% =====================================================================
%  Secao 10: Resolvendo as 4 Rainhas com Backtracking
% =====================================================================
\section{Resolvendo as 4 Rainhas com Backtracking}

O algoritmo inicia posicionando a primeira rainha. Em seguida, tenta posicionar a segunda. Caso encontre conflito, retorna e tenta outra posição. O processo continua até que todas as rainhas sejam posicionadas corretamente.

Para representar uma configuração do tabuleiro de forma compacta, usamos um vetor com quatro posições: cada posição corresponde a uma coluna e o número guardado indica em qual linha está a rainha daquela coluna. Assim, um único vetor descreve onde estão as quatro rainhas.

Uma solução válida é:
\[
  V = (3,\ 1,\ 4,\ 2)
\]

Isso significa:

\begin{itemize}[leftmargin=1.4em]
  \item Coluna 1 $\rightarrow$ Linha 3
  \item Coluna 2 $\rightarrow$ Linha 1
  \item Coluna 3 $\rightarrow$ Linha 4
  \item Coluna 4 $\rightarrow$ Linha 2
\end{itemize}

Observe que nenhuma rainha consegue atacar outra.

\figura{fig:tabuleiro-solucao}%
  {Solução válida do problema das 4 Rainhas: $V=(3,1,4,2)$.}%
  {figuras/tikz/tabuleiro}
